\documentclass{article} 

\usepackage[utf8]{inputenc} 
\usepackage[T1]{fontenc}    
\usepackage[english]{babel} 
\usepackage{amsmath}        
\usepackage{amssymb}        
\usepackage{amsfonts}       
\usepackage{graphicx}       
\usepackage{hyperref}       
\usepackage{geometry}       
\usepackage{enumitem}       
\usepackage{titling}
\usepackage{braket}
\usepackage{graphicx}
\usepackage{multicol}
\usepackage{lipsum}

% Page layout
\geometry{margin=1cm}      
\geometry{bottom=2cm}
\geometry{top=0.5cm}
\setlength{\parindent}{0pt} 
\setlength{\parskip}{6pt}   % Space between paragraphs
% math commands
\newcommand{\Lagrange}{\mathcal{L}}
\newcommand{\Hamilton}{\mathcal{H}}
\newcommand{\deriv}[3][]{\frac{d^{#1} #2}{d #3^{#1}}}
\newcommand{\pderiv}[3][]{\frac{\partial^{#1} #2}{\partial #3^{#1}}}

\renewcommand{\labelenumii}{\theenumii}
\renewcommand{\theenumii}{\theenumi.\arabic{enumii}.}

\graphicspath{{../plots}}

\pretitle{\begin{center}\LARGE}
\posttitle{\par\end{center}\vskip 0.5em}

\title{Metodo di Duhamel per equazione di ritenzione dei gas di fissione}
\author{Capitano A. B.}
\date{\today}

\begin{document}

\maketitle

\section{Metodo di Duhamel}
Sia \(u(r,t)\) la soluzione del problema di diffusione
$$\[
  \begin{cases}
    \pderiv{u}{t} = D \nabla^2 u + f(r,t)  \quad & r \in (0, R), t > 0 \\
    u(r,0) = 0 \quad & r \in [0, R], t = 0 \\
    + BC
  \end{cases}
\]

e \(w(r,t,s)\) famiglia di soluzioni con \(s \in [0,t]\) del problema omogeneo associato
\[
  \begin{cases}
    \pderiv{w}{t} = D \nabla^2 w  \quad & r \in (0, R), t > 0 \\
    w(r,0,s) = f(r,s) \quad & r \in [0, R], t = 0 \\
    + BC
  \end{cases} 
\]

per il principio di Duhamel la soluzione del problema non omogeneo è data da
\[
  u(r,t) = \int_0^t w(r,t-s,s) \, ds
\]

\section{Soluzione per gas di fissione}
\[
\begin{cases}
  \pderiv{u}{t} = D \frac{1}{r^2}\pderiv{}{r}( r^2 \pderiv{u}{r} ) + y\dot{F}  \quad & r \in (0, R), t > 0 \\
  u(r,0) = 0 \quad & r \in [0, R], t = 0 \\
  \pderiv{u}{r}(0,t) = 0 \quad & t > 0\\
  u(R, t) = 0 \quad & t > 0
\end{cases}
\]

dove \( u(r,t) \; [\text{atomi}/m^3]\) concentrazione di gas di prodotti dal rateo di fissione \(y\dot{F}\) e ritenuti in un grano di raggio R. Il problema omogeneo associato è
\[
\begin{cases}
  \pderiv{w}{t} = D \frac{1}{r^2}\pderiv{}{r}( r^2 \pderiv{w}{r} )  \quad & r \in (0, R), t > 0 \\
  w(r,0,s) = y\dot{F} \quad & r \in [0, R], t = 0 \\
  \pderiv{w}{r}(0,t) = 0 \quad & t > 0\\
  w(R, t, s) = 0 \quad & t > 0
  \end{cases}
\]

applicando una separazione di variabili \(w(r,t,s) = \theta(t,s)\phi(r)\)
\[
  \frac{1}{D\theta}\pderiv{\theta}{t} = \frac{1}{r^2\phi}\pderiv{}{r}( r^2 \pderiv{\phi}{r} ) = -k^2
\]

la soluzione della parte temporale è
\[
  \theta(t,s) = Ce^{-Dk^2t}
\]

espandendo la parte spaziale si ottiene
\begin{align*}
  \pderiv[2]{\phi}{r} + \frac{2}{r}\pderiv{\phi}{r} + k^2\phi = 0\\
  r\pderiv[2]{\phi}{r} + 2\pderiv{\phi}{r} + k^2r\phi = 0\\
  \pderiv[2]{}{r}(r\phi) + k^2r\phi = 0
\end{align*}
la cui soluzione è
\[
  \phi(r) = A\frac{\sin(kr)}{r} + B\frac{\cos(kr)}{r}
\]
applicando la condizione di Neumann omogenea in \(r = 0\) si ha \(B = 0\), mentre la condizione di Dirichlet in \(r = R\) impone
\[
  \sin(kR) = 0 \quad\Rightarrow\quad k_n = n\frac{\pi}{R}
\]
quindi la soluzione del problema omogeneo sarà della forma
\[
  w(r,t,s) = \sum_{n=1}^{\infty} b_n \frac{\sin(k_nr)}{r}e^{-Dk_n^2t}
\]
applicando ora la condizione iniziale si trovano i coefficienti di Fourier \(b_n\)
\[
  y\dot{F}r = \sum_{n=1}^{\infty} b_n \sin(k_nr)
\]
\[
  b_n = \frac{2y\dot{F}}{R}\underbrace{\int_0^R r\sin(k_nr)}_{R(-1)^{n+1} / k_n} \, dr = 2y\dot{F}\frac{(-1)^{n+1}}{k_n}
\]

si ha quindi la soluzione del problema omogeneo
\[
  w(r,t,s) = 2y\dot{F}\sum_{n=1}^{\infty} \frac{(-1)^{n+1}}{k_n}\frac{\sin(k_nr)}{r}e^{-Dk_n^2t} \quad [\text{atomi}/m^3s]
\]

Infine, integrando si ottiene la soluzione del problema non omogeneo
\[
  u(r,t) = \int_0^t w(r,t-s,s) \, ds = 2y\dot{F}\sum_{n=1}^{\infty} \frac{(-1)^{n+1}}{k_n}\frac{\sin(k_nr)}{r}\int_0^t e^{-Dk_n^2(t-s)} \, ds
\]
\[
  u(r,t) = \frac{2y\dot{F}}{D}\sum_{n=1}^{\infty} \frac{(-1)^{n+1}}{k_n^3}\frac{\sin(k_nr)}{r}\left(1 - e^{-Dk_n^2t}\right)\quad [\text{atomi}/m^3]
\]

Per trovare il numero di atomi ritenuti nel grano si integra la concentrazione su tutto il volume
\[
  N_R(t) = \int_V u(r,t) dV = 4\pi\int_0^R r^2 u(r,t) \, dr = \frac{8\pi y\dot{F}}{D}\sum_{n=1}^{\infty} \frac{(-1)^{n+1}}{k_n^3}\left(1 - e^{-Dk_n^2t}\right)\int_0^R r\sin(k_nr) \, dr
\]
esplicitando \(k_n\) nel risultato si ottiene
\[
  N_R(t) = \frac{8y\dot{F}R^5}{\pi^3D}\sum_{n=1}^\infty \frac{1 - e^{-D\left( n\pi/R \right)^2t}} {n^4}
\]

\end{document}

